%% toc: P1 8 8.4
\chapter{Motive Letters Instantiated by the Colon-Equals Sign}
\chaptermark{Instantiated Motive Letters}
\label{sec:p01-w02c-p01-018-motives}
\index{colon-equals sign!motive letter}\index{motive letter}\index{eliminator!motive}

\emph{Purpose.}  In the type-theoretic proofs of Part~12 the sign \(:=\)
often stands between a predicate letter applied to a variable and a
proposition, as in \(P(z):=(z=x)\).  Such a line can do two different
things.  It can introduce a new name for a predicate, which then keeps its
meaning to the end of the proof: role~(a) of
\cref{def:p01-w02b-p01-013-colon-equals}.  Or it can give the letter
that an eliminator or an induction principle binds in its statement a
value for one use of that principle: role~(b) of the same convention.
The two readings agree when the letter receives one value in a proof, and
they come apart when it receives several.  The proof
\cref{annex:p12.har.prime-equality} gives the letter \(P\) of the
equality eliminator three different values in one proof.  This chapter
states how the role of such a line is read, proves that the reading as
instantiation is the only one under which a letter with several values
has a single meaning at each of its occurrences, and records the role of
every line of this kind in the Part~12 proofs of the Harper development and
in the Part~7 proofs on higher inductive types, where the letter is a type
family with values in a universe.

\emph{Inputs.}  The colon-equals convention
\cref{def:p01-w02b-p01-013-colon-equals}, its role table
\cref{tab:p01-w02b-p01-013-roles} and the reading procedure
\cref{def:p01-w02b-p01-013-role-reading}; part~1 of
\cref{thm:p01-w02b-p01-013-colon-equals}, by which replacing a name of
role~(a) by its definiens never changes a truth value; the elimination
regime \cref{def:p12-har-prime-regime}, which states the induction
principle \(\mathsf{nat\_ind}\) with the bound letter \(P\) and the
propositional equality eliminator for predicates \(P:A\to\mathsf{Prop}\);
the lemmas \cref{thm:p12-har-prime-equality} and
\cref{thm:p12-har-prime-order}; and the statement of the generated
induction on terms in the autarkic part of the development
(\cref{def:p12-har-autarkic-syntax}); and the dependent induction
principles of the circle, the interval, the pushout and the set quotient
(\cref{def:p07-hit-circle}, \cref{def:p07-hit-interval},
\cref{def:p07-hit-pushout}, \cref{def:p07-hit-set-quotient}), each stated
for a type family \(P\) with values in a universe \(\mathcal U\).

\emph{Conclusions.}  \Cref{def:p01-w02c-p01-018-motive-reading} gives a
line \(p(\vec z):=s\) role~(b) when \(p\) is the letter bound by the
displayed statement of the rule applied at that place, and role~(a)
otherwise.  \Cref{thm:p01-w02c-p01-018-motive-readings} proves that the
two readings give the same meaning when the letter has one value in the
region, that no reading gives role~(a) to two such lines with the same
letter, and that a role-(a) reading of a letter with two different values
gives some occurrence two meanings.
\Cref{thm:p01-w02c-p01-018-equality-instances} computes the conclusions of
the three uses of the equality eliminator in
\cref{annex:p12.har.prime-equality} and shows that each of the two
region-wide role-(a) readings there either changes a conclusion or leaves
a term without a type.  \Cref{def:p01-w02c-p01-018-example} applies the
reading to every such line in the Part~12 proofs of the development, and
\cref{def:p01-w02c-p01-018-hit-example} to every such line in the Part~7
proofs on higher inductive types.

\section{Definitions and assumptions}
\label{sec:p01-w02c-p01-018-definitions}

The setting is a proof written in a typed calculus whose propositions form
a sort \(\mathsf{Prop}\), as in the elimination regime
\cref{def:p12-har-prime-regime}, or whose types form a universe
\(\mathcal U\), as in the higher inductive types of
\cref{sec:p07-higher-inductive-types}.  A \emph{region} is a region in the sense
of \cref{tab:p01-w02b-p01-013-roles}: a section, a body result or an annex
proof named in a row of that table.

\begin{hbdefinition}{def:p01-w02c-p01-018-motive}{Motive line, rule letter and application}
Let \(\rho\) be a rule, an eliminator or an induction principle whose
displayed statement binds a letter \(p\) of type
\(A_1\to\cdots\to A_k\to K\), where \(K\) is the sort \(\mathsf{Prop}\) or
a universe \(\mathcal U\), either by a product \(\Pi p\), or by a phrase
such as ``for \(P:S^1\to\mathcal U\)'' that fixes the letter before the
data of the principle, or
as a parameter written in the name of \(\rho\), as \(P\) is in
\(E_{A,x,P}\).  The predicate that \(\rho\) is used with is its
\emph{motive}.
\begin{enumerate}
  \item A \emph{motive line} in a region is an occurrence
  \(p(z_1,\ldots,z_k):=s\) with \(k\geq0\), also printed \(p\,z:=s\), in
  which
  \(z_1,\ldots,z_k\) are distinct variables of types
  \(A_1,\ldots,A_k\) and \(s\) is a term of \(K\) in the context extended
  by them: a proposition when \(K\) is \(\mathsf{Prop}\), a type of
  \(\mathcal U\) when \(K\) is \(\mathcal U\).  Its \emph{value} is
  \(\lambda z_1\ldots z_k.\,s\), called a \emph{predicate} in both cases.
  \item An \emph{application} of \(\rho\) is one use of \(\rho\) in the
  region: the stretch of the proof from the place where its motive is
  fixed to the conclusion that the use yields, including the premises
  that the use requires.
  \item The letter of a motive line is a \emph{rule letter at that place}
  if an application of a rule \(\rho\) whose displayed statement binds
  that same letter begins at the line, and the line fixes the motive of
  that application.  Otherwise the letter is a \emph{new name}.
  \item An occurrence of \(p\) in the region is \emph{bound by a
  statement} if it lies inside the displayed statement of a rule or
  result in which \(p\) is bound, or in a sentence that identifies two
  such statements.
\end{enumerate}
The \emph{meaning} of a free occurrence of \(p\) is the predicate it
stands for under a reading of the region.  Meanings are compared as
expressions up to renaming of bound variables: \(\lambda\vec z.\,s\) and
\(\lambda\vec z.\,s'\) are the same meaning exactly when \(s\) and \(s'\)
coincide after a renaming of \(\vec z\).
\end{hbdefinition}

Two assumptions are used.  First, the statement of the rule is displayed
before the motive line, in the region, in a declaration that the region
cites, or, for an annex proof, in the section that states the result it
proves; \cref{def:p12-har-prime-regime} is such a declaration for
\(\mathsf{nat\_ind}\) and the equality eliminator, and the definitions of
\cref{sec:p07-higher-inductive-types} are such statements for the induction
principles of the higher inductive types.  Second, applications
of rules in a region do not overlap: each motive line begins its own
application, and that application ends with its conclusion before the
next motive line with the same letter.  Both assumptions are checked for
each region in \cref{def:p01-w02c-p01-018-example}.

\section{The reading}
\label{sec:p01-w02c-p01-018-reading}

\begin{hbconvention}{def:p01-w02c-p01-018-motive-reading}{Role of a motive line}
A motive line \(p(\vec z):=s\) is read as follows.
\begin{enumerate}
  \item If its letter is a rule letter at that place, the line has
  role~(b) of \cref{def:p01-w02b-p01-013-colon-equals}, instantiation: the
  letter \(p\) of \(\rho\) takes the value \(\lambda\vec z.\,s\) for that
  one application, and each free occurrence of \(p\) inside the
  application means that value.  The same letter may take other values in
  other applications in the same region.
  \item Otherwise the line has role~(a), the definitional sign: \(p\) is a
  new name for \(\lambda\vec z.\,s\) from the line to the end of the
  region.  The left side must then be new in the region, so the region
  contains no second role-(a) line with the letter \(p\).
\end{enumerate}
\emph{Side conditions.}  The decision in item~1 uses the letter printed in
the statement of \(\rho\), not the letter's shape: a line \(Q\,n:=s\)
before an application of \(\mathsf{nat\_ind}\), whose statement binds
\(P\), has role~(a), and \(Q\) is a new name used as the motive.  A line
before an application of a rule whose statement binds no letter, as the
generated induction on formulas in \cref{thm:p12-har-prcompiler-form-typing}
is used there, has role~(a).  A region may contain lines of both roles; the
role table \cref{tab:p01-w02b-p01-013-roles} then lists the region in the
row of each role, with the lines named.
\end{hbconvention}

The convention is step~1 of \cref{def:p01-w02b-p01-013-role-reading}
applied to motive lines: a line whose left side is a letter of a rule
applied at that place has role~(b), and a line with a new name on its left
has role~(a).  It adds the test that decides which of the two holds for a
predicate letter applied to variables.

\section{Results}
\label{sec:p01-w02c-p01-018-results}

\begin{hblemma}{thm:p01-w02c-p01-018-motive-readings}{proved=annex:p01-w02c-p01-018-motive-readings}{The readings of a motive letter with one value and with several values}
Let \(R\) be a region and \(p\) a letter.  Let \(n\geq1\), and let
\(o_1,\ldots,o_n\) be the motive lines of \(R\) with the letter \(p\), in
this order,
with values \(\lambda\vec z.\,s_1,\ldots,\lambda\vec z.\,s_n\).  Suppose
that the letter of each \(o_i\) is a rule letter at its place, that no
other occurrence of \(:=\) in \(R\) has \(p\) on its left side, that the
applications \(\alpha_1,\ldots,\alpha_n\) begun by these lines are
pairwise disjoint, and that every occurrence of \(p\) in \(R\) lies in
some \(\alpha_i\) or is bound by a statement.
\begin{enumerate}
  \item If every \(o_i\) is read in role~(b), each free occurrence of
  \(p\) in \(R\) has exactly one meaning, namely
  \(\lambda\vec z.\,s_i\) for the occurrences in \(\alpha_i\).
  \item If \(n=1\), reading \(o_1\) in role~(a) gives every free
  occurrence of \(p\) in \(R\) the meaning \(\lambda\vec z.\,s_1\), the
  same as in item~1.
  \item If \(n\geq2\), no reading of \(R\) gives role~(a) to two of the
  lines \(o_1,\ldots,o_n\).
  \item If \(i<j\) and \(s_i\) and \(s_j\) are different propositions,
  even after renaming the variables \(\vec z\), then every reading that
  gives \(o_i\) role~(a) and \(o_j\) role~(b) gives the occurrence of
  \(p\) on the left of \(o_j\) two different meanings,
  \(\lambda\vec z.\,s_i\) and \(\lambda\vec z.\,s_j\).
\end{enumerate}
\end{hblemma}

Items~3 and~4 together say that when a letter receives two different
values in one region, a reading that gives any of its lines role~(a)
either breaks the newness condition of role~(a) or makes an occurrence
ambiguous; only the reading of item~1 gives each occurrence one meaning.

\begin{hblemma}{thm:p01-w02c-p01-018-equality-instances}{proved=annex:p01-w02c-p01-018-equality-instances}{Three uses of the equality eliminator in the Harper development}
Let \(A,B:\mathsf{Set}\), \(x,y,w:A\), \(f:A\to B\), \(e:x=y\), \(d:x=y\)
and \(e':y=w\), and let \(E_{A,a,P}(u,b,c):P\,b\) be the propositional
equality eliminator of \cref{def:p12-har-prime-regime}, for
\(P:A\to\mathsf{Prop}\), \(u:P\,a\) and \(c:a=b\).
\begin{enumerate}
  \item With \(P:=\lambda z.\,(z=x)\), the term
  \(E_{A,x,P}(\mathsf{refl\_equal}\ A\ x,\,y,\,e)\) has type \(y=x\).
  \item With \(P:=\lambda z.\,(f\,x=f\,z)\), where the equality is that of
  \(B\), the term
  \(E_{A,x,P}(\mathsf{refl\_equal}\ B\ (f\,x),\,y,\,e)\) has type
  \(f\,x=f\,y\).
  \item With \(P:=\lambda z.\,(x=z)\), the term
  \(E_{A,y,P}(d,\,w,\,e')\) has type \(x=w\).
  \item If \(P\) is read as one name for \(\lambda z.\,(x=z)\) in all three
  uses, the term of item~1 has type \(x=y\), the type of \(e\) itself.
  If \(P\) is read as one name for \(\lambda z.\,(z=x)\) in all three
  uses, the term of item~3 is well typed only if \(d\) has a type
  convertible with \(y=x\) under the relation \(=_{\beta\iota\delta}\) of
  \cref{def:p12-har-conversion}.  If the types of the variable \(d\) are
  exactly the types convertible with its declared type \(x=y\), if that
  relation is confluent, and if \(x\) and \(y\) are distinct variables,
  then the term has no type.
\end{enumerate}
\end{hblemma}

The lemma uses the base points that the uses in
\cref{annex:p12.har.prime-equality} need: \(x\) in the first two, and the
middle point \(y\) of the chain \(x=y\), \(y=w\) in the third, where the
seed is the given proof of \(x=y\).

\section{Worked example}
\label{sec:p01-w02c-p01-018-worked}

\begin{hbexample}{def:p01-w02c-p01-018-example}{Every motive line in the Part~12 proofs of the Harper development}
The proof \cref{annex:p12.har.prime-equality} opens by naming the
eliminator \(E_{A,x,P}(u,y,e):P\,y\) for \(u:P\,x\) and \(e:x=y\); its
letter \(P\) is bound by that statement and by the statement of
\(\mathsf{leib}\) in \cref{thm:p12-har-prime-equality}.  It then has three
motive lines.  ``Take \(P(z):=(z=x)\)'' begins a use of the eliminator
with seed \(\mathsf{refl\_equal}\ A\ x\), which yields
\(\mathsf{sym\_eq}\); ``take \(P(z):=(f x=f z)\)'' begins a use with the
reflexivity seed at \(f\,x\), which yields \(\mathsf{f\_equal}\); and
``the same transport with \(P(z):=(x=z)\)'' begins a third use, which
yields transitivity.  In each the letter is the rule letter \(P\) and the
line fixes the motive of the use that begins there, so by
\cref{def:p01-w02c-p01-018-motive-reading} all three have role~(b).  The
three uses end with their conclusions before the next line begins, and
the other occurrences of \(P\) in the proof are bound by the two
statements, so the hypotheses of
\cref{thm:p01-w02c-p01-018-motive-readings} hold with \(n=3\).  The values
\(z=x\), \(f\,x=f\,z\) and \(x=z\) are pairwise different propositions, so
by items~3 and~4 of that lemma the role-(b) reading is the only one that
gives each occurrence of \(P\) one meaning, and
\cref{thm:p01-w02c-p01-018-equality-instances} gives the conclusions
\(y=x\), \(f\,x=f\,y\) and \(x=w\).  The same proof also contains the line
\(\mathsf{no\_conf}\ n:=\lambda e.\,E(\mathsf{trivial},O,
\mathsf{sym\_eq}\ e)\), whose left side is a new name for a proof term;
it has role~(a).  The region is therefore listed in the rows of both
roles in \cref{tab:p01-w02b-p01-013-roles}, with the lines named.

Four other Part~12 proofs of the development each contain one motive line
whose letter is the rule letter.
\begin{itemize}
  \item \Cref{annex:p12.har.autarkic-square-root}: ``ordinary nat
  induction into the proposition \(P\,n:=\ldots\)''; the letter is the
  letter \(P\) of \(\mathsf{nat\_ind}\) in
  \cref{def:p12-har-prime-regime}.  The lines ``put \(M:=S\,n\) and
  \(q:=Q(S\,r)\)'' of the same proof introduce new names: role~(a).
  \item \Cref{annex:p12.har.autarkic-interpreter} and
  \cref{annex:p12.har.autarkic-substitution}: ``apply the generated
  \(T\) induction to \(P(t):=\ldots\)'' and ``use generated \(T\)
  induction into the precise proposition \(P(t):=\ldots\)''; the letter is
  the letter \(P\) of the generated induction on terms, stated for a fixed
  \(P:T\to\mathsf{Prop}\) in \cref{def:p12-har-autarkic-syntax}.
  The lines \(\nu_x:=\ldots\), \(X:=\mathsf{Var}\ O\) and
  \(\nu_\sigma:=\ldots\) introduce new names: role~(a).
  \item \Cref{annex:p12.har.prime-divisor}: ``set
  \(P\,n:=(\mathsf{one}<n)\to\mathsf{HPD}\ n\)'', followed by the use of
  \(\mathsf{cv\_ind}\), whose statement in
  \cref{thm:p12-har-prime-order} is made for every
  \(P:\mathsf{nat}\to\mathsf{Prop}\).  The line
  \(B_n:=\exists d.\,\mathsf{propdiv}\ d\ n\) introduces a new name, at
  which the classical principle \(\mathsf{tnd}\) is then used: role~(a).
\end{itemize}
In these four proofs the letter has one value, so by item~2 of
\cref{thm:p01-w02c-p01-018-motive-readings} the role-(a) reading of the
motive line would give every occurrence the same meaning; the role is
nevertheless~(b), because the letter is the rule letter.  The remaining
predicate lines of the development have role~(a): \(Q\,n:=\ldots\) in
\cref{annex:p12.har.prime-order}, whose letter \(Q\) is not the letter
\(P\) of \(\mathsf{nat\_ind}\) (there \(P\) already names the predicate of
the course-of-values theorem); \(P_0:=\lambda z{:}A.\,\mathrm{Eq}\,z\,a\)
in \cref{annex:p12.har.leibniz}, a new name for the predicate that is
then passed to a proof of Leibniz equality; and the correctness motive
\(P(p):=\ldots\) of \cref{thm:p12-har-prcompiler-form-typing} and
\cref{annex:p12.har.prcompiler-ok}, since the generated induction on
formulas is not displayed there with a bound letter.
\end{hbexample}

\begin{hbexample}{def:p01-w02c-p01-018-hit-example}{Every motive line in the Part~7 proofs on higher inductive types}
In \cref{sec:p07-higher-inductive-types} the induction principles of the
circle, the interval, the pushout and the set quotient are each stated for
a type family: ``for \(P:S^1\to\mathcal U\)'' in
\cref{def:p07-hit-circle}, ``into \(P:I\to\mathcal U\)'' in
\cref{def:p07-hit-interval}, ``for \(P:W\to\mathcal U\)'' in
\cref{def:p07-hit-pushout} and ``into a set-valued family
\(P:Q\to\mathcal U\)'' in \cref{def:p07-hit-set-quotient}.  The letter
\(P\) is therefore a rule letter, with \(K=\mathcal U\) in
\cref{def:p01-w02c-p01-018-motive}.  Five annex proofs each contain one
motive line with this letter, and in each the line fixes the family of the
one induction that begins there.
\begin{itemize}
  \item \Cref{annex:p07-hit-circle-recursion}: ``in \(P(x):=B\), supply
  \ldots\ to the dependent eliminator''; the constant family \(B\) is the
  value of the letter of \cref{def:p07-hit-circle}.
  \item \Cref{annex:p07-hit-circle-universal}: ``induct in
  \(P(x):=f(x)=g(x)\)''; the identity family is the value of the same
  letter.
  \item \Cref{annex:p07-hit-interval-contractible}: ``in
  \(P(x):=x=1_I\)''; the family of paths to \(1_I\) is the value of the
  letter of \cref{def:p07-hit-interval}.
  \item \Cref{annex:p07-hit-pushout-universal}: ``in
  \(P(x):=r(x)=t(x)\)''; the identity family is the value of the letter of
  \cref{def:p07-hit-pushout}.
  \item \Cref{annex:p07-hit-quotient-universal}: ``induct in the
  proposition \(P(x):=\|\sum_a(q(a)=x)\|\)''; the truncated fiber family is
  the value of the letter of \cref{def:p07-hit-set-quotient}, and the
  application ends with ``quotient induction gives all \(P(x)\)''.
\end{itemize}
In each proof the motive line is the only line with \(P\) on the left of
\(:=\), every other occurrence of \(P\) lies in the application it begins,
and so the hypotheses of \cref{thm:p01-w02c-p01-018-motive-readings} hold
with \(n=1\).  By item~2 of that lemma the role-(a) reading would give each
occurrence the same meaning as the role-(b) reading; by
\cref{def:p01-w02c-p01-018-motive-reading} the role is~(b), because the
letter is the letter that the displayed principle binds.  The other lines
of these annex files with \(:=\) introduce new names and have role~(a):
\(K(i):=\lambda a.k_a(i)\) in \cref{annex:p07-hit-interval-extensionality},
\(L(x):=\ldots\) in \cref{annex:p07-hit-hub-contraction}, and
\(t(x):=\ldots\) in \cref{annex:p07-hit-hub-dependent}.  The lines
\(H(r,r'):=u\) and \(S(r,r',x):=v(x)\) of the last proof are not motive
lines, because \(H\) and \(S\) are not of the type of a family; they give
the hub and spoke data letters \(H\) and \(S\) of the eliminator of
\cref{def:p07-hit-hub-truncation} their values for the one application
that follows (``apply the declared eliminator with \(g,H,S\)''), which is
role~(b) by step~1 of \cref{def:p01-w02b-p01-013-role-reading}.
\end{hbexample}

\section{Failure case}
\label{sec:p01-w02c-p01-018-failure}

\begin{hbwarning}{def:p01-w02c-p01-018-failure}{One name for three motives}
Reading the three motive lines of \cref{annex:p12.har.prime-equality} as
definitions of one name \(P\) fails in two ways.  The second and third
lines would redefine a name that the first has already introduced, which
role~(a) excludes (item~3 of
\cref{thm:p01-w02c-p01-018-motive-readings}).  And whichever value is kept
for the whole proof, some use of the eliminator proves something else or
nothing: with the last value \(\lambda z.\,(x=z)\) the first use
concludes \(x=y\), the type of the hypothesis \(e\) itself, instead of the
symmetric equation that \(\mathsf{sym\_eq}\) states; with the first value
\(\lambda z.\,(z=x)\) the seed of the third use, a proof of \(x=y\), must
have the type \(y=x\), and the transitivity term is well typed only if
\(d\) has a type convertible with \(y=x\), which item~4 of
\cref{thm:p01-w02c-p01-018-equality-instances} excludes for distinct
variables when the types of a variable are those convertible with its
declared type and conversion is confluent.  The handbook does not prove
these two properties for the conversion relation of
\cref{def:p12-har-conversion} in the historical system; the first failure,
the redefinition of a name, does not depend on them.  The
convention also does not decide a line whose rule is used without a
displayed statement binding the letter; such a line is a new name, and
\cref{def:p01-w02c-p01-018-motive-reading} then requires that no second
line in the region define the same letter.
\end{hbwarning}

\section{Relationships and sources}
\label{sec:p01-w02c-p01-018-relationships}

The convention \cref{def:p01-w02c-p01-018-motive-reading} refines
role~(b) of \cref{def:p01-w02b-p01-013-colon-equals} for predicate
letters, in the same way as that role covers the letters of the syllogism
derivations in \cref{annex:p02-w02-p02-002-twelve}, which also take
different values in different applications in one proof.  Item~2 of
\cref{thm:p01-w02c-p01-018-motive-readings} is the counterpart, for a
single motive line, of part~1 of \cref{thm:p01-w02b-p01-013-colon-equals}:
replacing a role-(a) name by its definiens changes nothing.  The rows
\(a:=b\) and \(p:=s\) of \cref{tab:p01-w02b-p01-013-roles} list the five
Part~12 proofs named in \cref{def:p01-w02c-p01-018-example} and the Part~7
proofs named in \cref{def:p01-w02c-p01-018-hit-example} in the row of each
role that their lines have, with the lines named.  The type-theoretic
background, products, \(\mathsf{Prop}\), and the eliminators generated by
inductive declarations, is that of the Harper development in
\cref{sec:p12-har-types}.  The first hypothesis of the last sentence of
item~4 of \cref{thm:p01-w02c-p01-018-equality-instances} holds in every
pure type system, whose conversion is \(\beta\)-conversion alone: by the
generation lemma, every type of a variable is \(\beta\)-convertible with
the type declared for it in the context
\cite[Lemma~5.2.13, item~2, p.~221]{Barendregt1992}.  The historical
system of \cref{def:p12-har-prime-regime} adds inductive types and their
computation rules to such a system, and the handbook proves neither that
hypothesis nor confluence for the extended conversion relation.

\emph{Sources.}  The motive lines and the statements of the rules are
quoted from the Part~12 presentation of the historical arithmetic library
and Euclid development, \cref{def:p12-har-prime-regime},
\cref{thm:p12-har-prime-equality}, \cref{thm:p12-har-prime-order} and
their annex proofs, and from its autarkic part,
\cref{def:p12-har-autarkic-syntax}, \cref{thm:p12-har-autarkic-interpreter},
\cref{thm:p12-har-autarkic-substitution} and their annex proofs; the
bibliographic sources of that presentation are cited there.  The Part~7
lines and induction principles are quoted from
\cref{sec:p07-higher-inductive-types} and its annex proofs, where their
sources are cited.  The convention, the two lemmas and
their proofs are editorial.
